\section{Results for the log calibration}\label{sec:main}

Throughout this section $P$ is a household economy with $\underline a=0$, cap $\bar a$, log
utility, $\beta=24/25$, and an earnings process $(Z,y,\Pi)$. Two conditions on the earnings
process are assumed in every uniqueness statement:
\begin{enumerate}[label=(E\arabic*)]
  \item\label{E1} the transitions are monotone;
  \item\label{E2} there is a state $z_0$ with $y(z_0)=y_{\min}$ and $\Pi(z,z_0)>0$ for every $z$.
\end{enumerate}
The cap enters through one inequality, evaluated at the larger of the two rates compared:
\begin{equation}\label{eq:cap}
  2\beta\,y_{\max}<\bigl(1-\beta(1+r)\bigr)\,\bar a .
\end{equation}
It says the cap was set large enough not to bind; nothing else about $\bar a$ is used.

\subsection{Uniqueness}

\begin{theorem}[Aiyagari at $\mu=1$: the equilibrium rate is unique on $(-\delta,\lambda)$]
\label{thm:main}
Let $r_1,r_2\in(-2/25,\,1/24)$ be admissible, and suppose that for $i=1,2$ there is a stationary
distribution $\mu_i$ of $P.\text{withRate}(r_i)$ with $K(\mu_i)=D(r_i)$, where
$D(r)=\alpha/((1-\alpha)(r+\delta))$ with $\alpha=9/25$, $\delta=2/25$. If \eqref{eq:cap} holds at
$\max\{r_1,r_2\}$, then $r_1=r_2$.
\formal{IncomeFluctuation.aiyagari1994\_log\_equilibriumRate\_unique\_Ioo}
\end{theorem}

The theorem is stated for every interval $[r_{lo},r_{hi}]$ with $-\delta<r_{lo}\le r_{hi}<\lambda$
in the form \lean{aiyagari1994\_log\_equilibriumRate\_unique}, with \eqref{eq:cap} at $r_{hi}$;
the open-interval version applies it to $[\min\{r_1,r_2\},\max\{r_1,r_2\}]$. The lower end
$-\delta$ is where demand becomes infinite; the upper end $\lambda$ is where the household becomes
patient. Both are excluded for a reason (Theorems~\ref{thm:above} and \ref{thm:below}), so
Theorem~\ref{thm:main} covers every rate at which an equilibrium could exist, up to a band above
$\lambda$ discussed below.

\begin{corollary}[Aiyagari's own earnings processes]\label{cor:tauchen}
Let $n\ge0$, $\rho\in[0,1)$, $\sigma>0$, $c>0$, and let the endowment take the $n+2$ values
$c\,e^{-3\sigma+i\cdot 6\sigma/(n+1)}$, $i=0,\dots,n+1$, with the Tauchen transition matrix of
$\log l'=\rho\log l+\sigma\sqrt{1-\rho^2}\,\varepsilon$ on that grid. Then the conclusion of
Theorem~\ref{thm:main} holds, with \eqref{eq:cap} read with $y_{\max}=c\,e^{3\sigma}$.
\formal{aiyagariTauchen\_equilibriumRate\_unique\_Ioo}
\end{corollary}

The corollary covers all eight of Aiyagari's processes and any finer grid. Its proof is that the
Tauchen chain satisfies \ref{E1} and \ref{E2}: the cumulative rows
$F_i(k)=\Phi\bigl((x_k+h/2-\rho x_i)/\sigma_\varepsilon\bigr)$ are decreasing in $i$ because $\Phi$
is monotone and $\rho\ge0$, which is first-order dominance, and every entry of the first column is
a value of $\Phi$ on the real line, hence positive. The normal distribution is Mathlib's
$\mathrm{gaussianReal}$, so $\Phi$ is a genuine cumulative distribution function and the
positivity is a theorem about the Lebesgue integral, not a numerical fact. The scale $c$ is
whatever normalises the mean to one; uniqueness does not depend on it except through
\eqref{eq:cap}.

\begin{remark}[What is used about the earnings process]
Only \ref{E1} and \ref{E2}. No transition probability is named, no stationary distribution of the
chain is computed, and the number of states is arbitrary. \ref{E1} is what Light's Theorem 1 needs
to propagate monotonicity of the value function in income through the expectation; \ref{E2} is
what the Doeblin argument needs to produce an atom. For $\rho=0$ both are immediate; for $\rho>0$
they are Huggett's Lemma 1 for Tauchen chains.
\end{remark}

\subsection{No equilibrium outside \texorpdfstring{$(-\delta,\lambda)$}{the interval}}

\begin{theorem}[Above the rate of time preference]\label{thm:above}
Let $r\ge0$ with $\beta(1+r)>1$, and let $\mu$ be stationary for $P.\text{withRate}(r)$. If
\[
  y_{\max}+(1+r)D<\Bigl(1-\frac{1}{\beta(1+r)}\Bigr)\bigl(y_{\min}+r\bar a\bigr),
\]
then $K(\mu)\ne D$. In particular no equilibrium exists at $r$ once $r-\lambda$ exceeds a
quantity of order $(y_{\max}+(1+r)D(r))/(\beta r\bar a)$.
\formal{IncomeFluctuation.log\_no\_equilibrium\_of\_patient\_risk}
\end{theorem}

The proof is a supermartingale bound. With $\beta R>1$ the Euler inequality makes
$u'(c_t)$ a supermartingale that cannot decrease on average, so consumption cannot rise on average,
while cash on hand grows at rate $r$ on the part of the distribution near the cap; integrating
against the stationary distribution gives a floor on $K(\mu)$ of order $(1-1/\beta R)\,r\bar a$.
The band $(\lambda,\lambda+\epsilon)$ that escapes has width shrinking like $1/\bar a$, and with
$\bar a$ chosen to satisfy \eqref{eq:cap} it is thin. A version for CRRA with integer risk
aversion $n$, via Jensen on $x^{-n}$, is \lean{crra\_no\_equilibrium\_of\_patient\_mean}.

\begin{theorem}[Negative rates]\label{thm:below}
Assume \ref{E2} and that the earnings process is normalised so that the invariant distribution of
$\Pi$ has mean at most one. Let $\mu$ be stationary for $P.\text{withRate}(r)$.
\begin{enumerate}[label=(\roman*)]
\item For any period utility satisfying the standing assumptions, if $r\in(-2/25,-13/250]$ then
  $K(\mu)\ne D(r)$.
\item For log utility, if $24\,y_{\max}<\bar a$ and $r\in(-2/25,-1/25]$ then $K(\mu)\ne D(r)$.
\end{enumerate}
\formal{IncomeFluctuation.aiyagari1994\_no\_equilibrium\_of\_budget; IncomeFluctuation.aiyagari1994\_log\_no\_equilibrium\_below\_mean}
\end{theorem}

Part (i) uses nothing about preferences: at a negative rate, integrating the budget constraint
against a stationary distribution gives $K\le \bar Y/(-r)$, the capital whose interest cost
exhausts mean income, which is the zero-consumption supply schedule of \citet{YoungWalsh2025}.
Part (ii) uses the log ceiling $K\le\beta\bar Y/(1-\beta(1+r))$ from the minimal MPC
(Section~\ref{sec:mpc}). Both need the income marginal of the stationary distribution to be the
chain's invariant distribution, which is where \ref{E2} enters, through a finite-chain uniqueness
argument (\lean{invariant\_eq\_of\_reach}). In both cases the ceiling falls below $D(r)$, which
grows without bound as $r\downarrow-\delta$.

Together, Theorems~\ref{thm:main}--\ref{thm:below}: for Aiyagari's log calibration the set of
equilibrium rates is contained in $(-1/25,\lambda+\epsilon)$ and has at most one point in
$(-\delta,\lambda)$.

\subsection{Existence}\label{sec:exist}

\begin{theorem}[Existence on an interval, given a supply floor at its top]\label{thm:exist}
Let $-\delta<r_{lo}<r_{hi}<\lambda$, assume \ref{E1}, \ref{E2} and \eqref{eq:cap} at $r_{hi}$.
Suppose
\[
  \frac{\beta\,y_{\max}}{1-\beta(1+r_{lo})}\le D(r_{lo}),\qquad
  D(r_{hi})\le m\le K(\mu)\ \text{ for every stationary }\mu\text{ at }r_{hi}.
\]
Then there is an equilibrium rate in $[r_{lo},r_{hi}]$.
\formal{IncomeFluctuation.aiyagari1994\_log\_exists\_equilibrium\_of\_floor}
\end{theorem}

The first inequality is the log capital ceiling at $r_{lo}$ against demand there, and it holds
for $r_{lo}$ near $-\delta$ because demand is unbounded. The second is a floor $m$ on capital
supply at $r_{hi}$ that exceeds demand there, and it is the one quantitative fact about the log
household that this development does not prove. Aiyagari's equilibria sit where $k/w\approx5$; the
analytic supply floors available (an atom at the cap, positive saving from the corner) are nowhere
near that, and how much an impatient household accumulates before impatience takes over is a
quantitative question about the policy that the qualitative theory does not answer. With the floor
supplied, existence follows by the intermediate value theorem, because capital supply along the
unique stationary distribution is continuous in the rate (Section~\ref{sec:cont}).

\subsection{Computation}

\begin{theorem}[Bisection converges to the equilibrium rate]\label{thm:bisect}
Under the hypotheses of Theorem~\ref{thm:main} on $[r_{lo},r_{hi}]$, let $\nu$ be any selection of
stationary distributions, $\nu(r)$ stationary for $P.\text{withRate}(r)$ on $[r_{lo},r_{hi}]$,
and let $E(r)=K(\nu(r))-D(r)$ be excess supply. If $r^*\in[r_{lo},r_{hi}]$ is an equilibrium rate,
then the bisection on $E$ started from $[r_{lo},r_{hi}]$ brackets $r^*$ at every step, its
midpoint after $n$ steps is within $(r_{hi}-r_{lo})/2^{n+1}$ of $r^*$, and its endpoints converge
to $r^*$.
\formal{IncomeFluctuation.aiyagari1994\_log\_bisection}
\end{theorem}

Excess supply is strictly increasing under these hypotheses, because supply is monotone
(Section~\ref{sec:coupling}) and demand strictly decreasing, and strict monotonicity alone decides
which half of a bracket contains the zero; neither continuity nor a sign condition at the ends is
used. Any stationary distribution at $r$ is the selection's, by the Doeblin uniqueness, so the
algorithm does not depend on how the stationary distribution is computed. The theorem certifies
the outer loop of the standard procedure, as in \citet{KirkbyToolkit2017}, given exact inner
solutions; convergence of the discretised inner solutions to the true ones is the subject of
\citet{Kirkby2017} and \citet{Kirkby2019}, and is not re-proved here.
